% !TEX root = ../main.tex

\chapter{Conclusions}

\section{Discussions}

This thesis demonstrates how a mobile literature assistant can preserve continuity across changing papers, generated assistance, and evolving research interests.
The completed system maintains explicit identities for papers, revisions, evidence, behavioral events, and mobile presentation states, and the evaluations verify these mechanisms under documented configurations.
The experiments also identify the empirical boundaries that guide further validation.

RQ1 concerned revision-correct paper identity and section-level provenance.
The data model separates a logical paper from each observed revision and propagates revision identity, content hashes, and transformation configuration to derived artifacts.
Deterministic preparation and idempotent jobs make repeated processing reproducible, while revision-scoped retrieval prevents a question about one revision from drawing chunks from another.
The corresponding software tests establish the implemented isolation property.
They leave parser structure quality, stale-artifact invalidation, and recovery performance across diverse scientific layouts as open empirical questions.

RQ2 asked how retrieval and generation can remain inspectable against paper evidence.
The implemented pipeline restricts retrieval to one paper revision, combines dense similarity with a fixed lexical reranker, refuses generation below an evidence threshold, validates citation markers, and reports a bounded source-match state.
In the fifteen-case live evaluation, all answered cases contained a verified citation and passed the lexical source-match rule.
None of the twelve answerable questions was falsely refused, while two of the three unanswerable cases followed the specified refusal-marker protocol.
Dense retrieval reached the designated section in 4 of 12 answerable cases.
This gap shows that a supported answer may draw on a different section from the passage selected by the case author, particularly when central claims recur in abstracts and introductions.
The current evaluation reports both outcomes, preserving the weaker section-placement result alongside successful source matching.

The model-tier comparison adds an operational result to RQ2.
The flagship and mid-tier routes both produced citations for every answered case, but the mid-tier route emitted no required refusal marker on the three unanswerable cases in the reported run.
The same mid tier satisfied all measured structural criteria for seven paper summaries, with mean lexical claim support of 0.95 compared with 0.94 for the flagship tier, while reducing measured cost from \$0.921 to \$0.140.
These results support task-specific routing under the measured criteria: the mid tier is suitable for structurally checked summarization, whereas the current evidence favors the flagship route for grounded question answering.
Human comparison is still required to evaluate explanatory quality, and lexical source matching cannot establish semantic entailment or scientific truth.

RQ3 concerned adaptation from explicit interests and bounded behavioral evidence.
The proposed model uses exposure-qualified negative feedback, separate positive and negative centroids, temporal decay, per-paper saturation, evidence confidence, and deterministic recommendation reasons.
Across eleven applicable controlled traces, its mean nDCG@5 of 0.92 was close to the recency-only baseline of 0.92.
Its reciprocal rank was 0.94, compared with 0.88 for the recency-only and earlier single-timescale baselines.
Mechanism-level comparisons showed the intended response to recent positive evidence, exposed negative feedback, unexposed skips, repeated interaction with one paper, and duplicate replay.
Per-paper saturation increased nDCG@5 from 0.51 to 0.73 in the repeated-paper trace, and duplicate replay produced no score change.
The separate explicit-correction experiment exercised topic addition, replacement, and deletion with eight equal-age candidates. All four directional checks passed: adding a topic changed its score and reason, while replacement and deletion changed the affected topic group's mean rank by two positions and removed obsolete match reasons.
The scenario set supports the mechanics and failure protections of the model.
It produced no independent aggregate ranking gain for dual-timescale decay or confidence gating; real-user preference accuracy and long-term benefit remain questions for longitudinal evaluation.

RQ4 addressed whether the integrated mobile workflow preserves paper identity, evidence access, and interaction meaning while exposing live, cached, processing, and failure states.
The Android baseline completed all 810 controlled measurements of startup, state transitions, polling, graph interaction, and event synchronization.
The resource experiment added 12{,}960 successful samples across sixteen independent cold-boot sessions.
In the live integration trace, all 25 event stages met their conditions, including retry after an injected failure, duplicate-aware ingestion, and client readback.
The complete device path then produced a five-paper briefing, opened paper detail, returned a source-linked answer, rendered a multi-node citation graph, and opened a selected neighboring paper.
Two further arXiv seeds completed the seed-to-graph path. These repetitions provide limited evidence beyond the original seed, rather than a general claim about all papers.

The mobile measurements also identify where time is spent.
In the baseline Android~14 configuration, cold process start had a median of 2.77~s, whereas foreground/resume had a median of 134~ms.
Updating an existing graph WebView remained below 112~ms at the tested sizes, while creating a new WebView was slower and more variable.
The resource-controlled study indicates that startup was more sensitive to assigned CPU and memory than existing-WebView updates or native graph selection.
Its new-WebView timer also developed a synchronization-sensitive five-second mode under the four-core settings. Those samples remain in the record, but they are excluded from causal interpretation of the CPU effect.
The live traces show the larger system boundary: generation-inclusive briefing and answer stages took 22.4~s and 24.3~s, compared with 4.7~s and 0.9~s after backend artifacts could be reused.
Because these are separate single observations rather than repeated cold and warm trials, they locate likely waiting time but do not estimate a general speed-up.

The formative prototype study complements the software measurements.
All five participants completed the briefing, source, and cited-Q\&A tasks; four completed interest editing, and two completed graph exploration.
The graph result led to a visible instruction, persistent selected-node state, relation context, and a stable \emph{Open Paper} action.
Software tests cover these changes, and participant retesting remains future work.
The five-person sample, moderated prompts, fixed prototype content, and absence of participant-level timing and interaction records limit the study to design-stage observations.

The combined evidence supports four contributions.
The revision-aware data and job model makes paper-derived artifacts attributable to a defined source state.
The assistance pipeline gives retrieved evidence, refusal, and source matching explicit interfaces and failure outcomes.
The behavioral model makes event semantics, confidence, and recommendation reasons replayable and inspectable.
The Android client carries these properties through cached state, asynchronous polling, duplicate-safe event delivery, and bounded graph interaction.
The evidence comes from complementary forms of testing rather than one aggregate score: backend regression tests establish software contracts, controlled traces isolate mechanisms, live question-answering cases exercise the provider boundary, and Android measurements cover the client state machine and integration path.
Failed or anomalous outcomes are retained when they change the interpretation, including the incomplete refusal result and the unstable new-WebView timing mode.

The completed evaluation defines a clear scope for the conclusions.
The document-processing evaluation does not cover a stratified corpus of scientific layouts and revisions.
The question set contains fifteen manually curated cases, and its source-match metric is lexical.
The behavior scenarios are designed tests with hand-specified relevance judgments.
The client measurements use one Android version, device profile, and emulator host, and the usability study contains five participants.
Longitudinal recommendation benefit, representative answer quality, physical-device variability, and research-productivity effects define the next stages of evaluation.

\section{Main Conclusions}

MNEME addresses the work that follows paper discovery: deciding what deserves attention, preserving what has been read, inspecting machine-generated explanations, and reconnecting related papers over time.
Its central method is a continuous mobile reading loop whose paper identities, evidence references, behavioral signals, and visible states share explicit contracts.

The completed evaluation shows that the implemented mechanisms operate together in the recorded environment.
Revision-scoped retrieval preserves the tested paper-version boundary.
The fifteen-case question evaluation produced cited, source-matched answers for every answered case, avoided false refusal on all twelve answerable cases, and followed the specified refusal protocol in two of three unanswerable cases.
The controlled behavior evaluation confirmed the intended effects of exposure gating, negative feedback, per-paper saturation, and duplicate-safe replay; reciprocal rank improved while aggregate nDCG remained close to the recency baseline.
The Android experiments demonstrated cache restoration, asynchronous polling, multi-node graph interaction, duplicate-safe event synchronization, and a live seed-to-reading path.

The evidence supports an implementation-grounded system contribution with bounded empirical results.
General claims about research productivity, long-term recommendation quality, and scientific answer correctness require broader longitudinal and human evaluation.
The main lesson is that an AI literature assistant needs inspectable boundaries around paper versions, retrieved evidence, inferred interests, and degraded mobile state.
Those boundaries make every output state available for systematic evaluation and improvement.

\section{Outlook}

Document processing should be evaluated on a stratified collection of scientific PDFs covering multiple disciplines, layouts, revision histories, and degraded files.
The study should measure section recovery, chunk attribution, invalidation after revision, failure recovery, and processing latency.

Evidence-grounded assistance requires a larger question set with human claim--evidence judgments.
Future evaluation should compare retrieval and verification ablations, calibrate refusal thresholds, and measure semantic support in addition to lexical source matching.
The cheaper model route should be reconsidered for question answering after its refusal behavior is corrected and retested.

Interest adaptation requires longitudinal evidence.
A study with real reading histories should compare explicit-only, recency-based, and behavioral models while measuring relevance, diversity, stability, and the effect of user correction.
The hand-specified weights and time constants should be treated as hypotheses for that study.

Mobile evaluation should add physical devices with different performance and memory profiles, retain participant-level interaction records, and repeat the graph task after the implemented redesign.
Accessibility testing, notification behavior across Android versions, privacy controls for deletion and export, and longer offline periods also require direct evaluation.
Concept and author graph layers can be examined after the bounded paper graph has been validated with users.

\clearpage
